

%\begin{definition}[Diagrams]
%	Let $\category I$ and $\category C$ be categories. A \emph{$\category I$-shaped diagram} in the category $\category C$ is a functor $\category I \to \category C$.
%\end{definition}

%\begin{definition}[Constant Diagrams]
%	Let $\category I$ and $\category C$ be categories. For $X \in \category C$ the \emph{constant diagram} at $X$ is the constant functor $\ConstantFunctor X\colon \category I\to\category C$ given by $i\mapsto X$ and $(i\to j)\mapsto \id_X$.
%\end{definition}


%\begin{definition}[Cones]
%	Let $F\colon \category I\to \category C$ be a diagram and $X\in\category C$. A \emph{cone over} $F$ with \emph{tip} $X$ is a natural transformation $\ConstantFunctor X \Rightarrow F$. Dually, a \emph{cocone under $F$} with \emph{tip} $X$ is a natural transformation $F\Rightarrow \ConstantFunctor X$.
%\end{definition}

%\begin{definition}[Universality]
%	Write down cones explicitly? I don't like this definition. Its horribly inexplicit what a limit should be.
%	Let $F\colon \category I\to\category C$ be a diagram. A cone $c\colon\ConstantFunctor{L} \Rightarrow F$ is called \emph{universal} if for any cone $c'\colon\ConstantFunctor{L'} \Rightarrow F$ there is a unique morphism $u\colon L'\to L$ such that $c\circ \hat u = c'$ where $\hat u\colon \ConstantFunctor{L'} \Rightarrow \ConstantFunctor{L}$ is such that $\hat u_i = u$ for all $i\in \category I$. Dually, a cocone $c\colon F\Rightarrow \ConstantFunctor{C}$ is called \emph{universal} if for any cocone $c'\colon F\Rightarrow \ConstantFunctor{C'}$ there is a unique morphism $u\colon C \to C'$ such that $\hat u \circ c = c'$ where again $\hat u\colon \ConstantFunctor{L'} \Rightarrow \ConstantFunctor{L}$ is such that $\hat u_i = u$ for all $i\in \category I$.
%\end{definition}

%\begin{definition}[Limits]
%	Let $D\colon \category I\to\category C$ be a diagram. A universal (co)cone with tip $T$ is called a \emph{(co)limit} of $D$. We write 
%		$$\limit D \ldef \limit_\category I D\ldef \limit_{i\in\category I} D(i) \ldef T\text{ and } \colimit D \ldef \colimit_\category I D\ldef \colimit_{i\in \category I} D(i) \ldef T$$
%	for the tip of the limit and colimit respectively.
%\end{definition}

\begin{definition}[Reflexive pairs and reflexive (co)equalizers]
	A parallel pair of morphisms $X\rightrightarrows Y$ with a common section $Y\to X$ is called a \emph{reflexive pair}. A (co)limit of a diagram with underlying shape a reflexive pair is called a \emph{reflexive (co)equalizer}.
\end{definition}

\begin{definition}[Filteredness]
	Let $\kappa$ be an infinite regular cardinal. A category $\category I$ is called \emph{$\kappa$-filtered} if any $\kappa$-small diagram in $\category I$ admits a cocone. In the case of $\kappa=\N$, we simply call $\category I$ \emph{filtered}. Dually, $\category I$ is called \emph{$\kappa$-cofiltered} or \emph{cofiltered} if $\category I^\op$ is $\kappa$-filtered or filtered respectively.
\end{definition}

\begin{remark}[Interaction of filtered colimits and limits]
	\label{remark: interaction of filtered colimits and limits}
	Let $\kappa$ be an infinite regular cardinal. Then $\kappa$-filtered colimits of sets commute with $\kappa$-small limits. Explicitly, if $\category I$ is $\kappa$-filtered, $\category J$ is $\kappa$-small and $D\colon \category I\times\category J \to \Set$ is a diagram, then the canonical morphism
		$$\colimit_\category I \limit_\category J D \to \limit_\category J \colimit_\category I D$$
	is an isomorphism. The author could not locate a reference for this claim except for an $\infty$-categorical analogue given by Lurie in \cite[Prop. 5.3.3.3]{higher-topoi}. This is however no obstruction as one can represent any ordinary category by its nerve, apply the theorem and recover the categories as the underlying homotopy categories of the nerves. This works since (homotopy) limits and colimits in the nerve agree with the limits and colimits in the ordinary category.
\end{remark}

%\begin{definition}[Siftedness]
%	\label{definition: siftedness}
%	
%	A category $\category I$ is called \emph{sifted} if colimits of $\Set$-valued diagrams of shape $\category I$ commute with finite products. Explicitly for every diagram $D\colon \category I \times \category S \to \Set$ where $\category S$ is finite and discrete (\ie the only morphisms are the various identities), the canonical morphism
%		$$\colimit_{\category I} \prod_{\category S} D \to \prod_{\category S} \colimit_{\category I} D$$
%	is an isomorphism.
%	
%	\begin{todo}
%		Show that filtered colimits and reflexive coequalizers are sifted
%	\end{todo}
%\end{definition}

\begin{construction}[Pro-categories]
	\label{Construction: Pro-categories}
	
	For a category $\category C$ the category $\ProCategory {\category C}$ of \emph{pro-objects} of $\category C$ essentially arises from formally adjoining all cofiltered limits missing in $\category C$.
	
	More precisely, consider the fully faithful \emph{Yoneda embedding}
	\begin{align*}
		\yoneda\colon \category C^\op &\to \functorCategory {\category C} {\Set},\\
		X&\mapsto \Hom_{\category C^\op}(-, X) = \Hom_\category C(X, -),\\
		Y\xrightarrow{f^\op} X &\mapsto (f^\op)_\ast = f^\ast \colon \Hom_\category C(Y, -)\to \Hom_\category C(X, -)
	\end{align*}
	of $\category C^\op$. Consider its opposite $\yoneda^\op\colon \category C \to \functorCategory{\category C}{\Set}^\op$. Now $\ProCategory{\category C}$ is the full subcategory of $\functorCategory{\category C}{\Set}^\op$ of cofiltered limits of representable functors, \ie
		$$\ProCategory{\category C} \ldef \left\{\limit \Hom_\category C(X, -)\setseparator X\colon \category I\to\category C\text{ cofiltered diagram}\right\}\subseteq_\text{full} \functorCategory{\category C}{\Set}^\op$$
	where $\Hom_\category C(X, -)$ is to be understood as $\yoneda^\op \circ X$.
\end{construction}

%\begin{remark}
%	In the setting of construction \ref{Construction: Pro-categories}, let $F = \limit_{i\in\category I} \Hom_\category C(X_i, -)$ and $F'=\limit_{j \in \category J} \Hom_\category C(X_j', -)$ be pro-objects of $\category C$, such that $X$ and $X'$ are cofiltered. Then for any $Y\in\category C$ we have a natural isomorphism
%		$$F(Y) \cong  \limit_{i\in \category I} \Hom_\category C(X_i, Y).$$
%	Similarly there is a natural isomorphism
%		$$\Hom_\ProCategory{\category C}(F, F') \cong \limit_{j\in \category J}\colimit_{i\in \category I} \Hom_\category C(X_i, X_j').$$
%\end{remark}

%\begin{lemma}[Pro-categories as subcategories]
%https://ncatlab.org/nlab/show/pro-object, 3.1
%	\begin{todo}
%		Define cocompact
%		
%		Should I use $\subseteq_full$, should I define it, or should I leave it out completely?
%	\end{todo}
%	Let $\category C\subseteq_\text{full} \category C'$ be a full subcategory and assume that $\category C'$ has all cofiltered limits and that each $X\in \category C$ is cocompact in $\category C'$. Then there is a natural full and faithful functor
%	\begin{align*}
%		\ProCategory{\category C} &\hookrightarrow \category C',\\
%		\limit \Hom_\category C(X, -) &\mapsto \limit X
%	\end{align*}
%\end{lemma}
%
%\begin{proof}
%	\begin{todo}
%	\end{todo}
%\end{proof}
